\documentclass[11pt,a4paper]{article}
\usepackage[utf8]{inputenc}
\usepackage[T1]{fontenc}
\usepackage{amsmath,amssymb,amsthm}
\usepackage{geometry}
\geometry{margin=2.5cm}
\theoremstyle{plain}
\newtheorem{theorem}{Theorem}[section]
\newtheorem{proposition}[theorem]{Proposition}
\newtheorem{lemma}[theorem]{Lemma}
\theoremstyle{definition}
\newtheorem{definition}[theorem]{Definition}
\title{Soluciones Tests 16-20: Álgebra}
\author{Mathematical AI Agent}
\date{\today}
\begin{document}
\maketitle
\section{Problemas}

\subsection*{Test 16: Pequeño teorema de Fermat}

\begin{theorem}[Pequeño teorema de Fermat]
Sea $p$ un número primo y $a \in \mathbb{Z}$. Entonces $a^p \equiv a \pmod{p}$.
\end{theorem}

\begin{proof}
Procedemos por inducción sobre $a$.

\textbf{Caso base:} Para $a = 0$, tenemos $0^p = 0 \equiv 0 \pmod{p}$, lo cual es verdadero.

\textbf{Paso inductivo:} Supongamos que el resultado vale para un entero $a \geq 0$, es decir, $a^p \equiv a \pmod{p}$. Debemos probar que $(a+1)^p \equiv a+1 \pmod{p}$.

Por el binomio de Newton:
\[
(a+1)^p = \sum_{k=0}^{p} \binom{p}{k} a^k = a^p + \sum_{k=1}^{p-1} \binom{p}{k} a^k + 1.
\]

Para $1 \leq k \leq p-1$, el coeficiente binomial $\binom{p}{k} = \frac{p!}{k!(p-k)!}$ es divisible por $p$, ya que $p$ es primo y aparece en el numerador pero no en el denominador (pues $k < p$ y $p-k < p$). Por lo tanto, $\binom{p}{k} \equiv 0 \pmod{p}$ para $1 \leq k \leq p-1$.

Así:
\[
(a+1)^p \equiv a^p + 1 \pmod{p}.
\]

Por hipótesis inductiva, $a^p \equiv a \pmod{p}$, de donde:
\[
(a+1)^p \equiv a + 1 \pmod{p}.
\]

Para el caso $a < 0$, basta notar que $(-a)^p \equiv -a \pmod{p}$ por lo demostrado, y como $p$ es primo impar (salvo $p=2$ que se verifica directamente), tenemos $(-a)^p = -a^p$, luego $-a^p \equiv -a \pmod{p}$, es decir, $a^p \equiv a \pmod{p}$.
\end{proof}

\subsection*{Test 17: Infinitud de primos}

\begin{theorem}[Euclides]
La conjunto de números primos es infinito.
\end{theorem}

\begin{proof}
Supongamos, por contradicción, que existen solo un número finito de primos, digamos $p_1, p_2, \ldots, p_n$. Consideremos el número:
\[
N = p_1 p_2 \cdots p_n + 1.
\]

Claramente $N > 1$, por lo que $N$ tiene al menos un divisor primo, digamos $q$.

Si $q$ fuera igual a algún $p_i$ con $1 \leq i \leq n$, entonces $q$ dividiría al producto $p_1 p_2 \cdots p_n$. Pero como $q$ divide a $N$, tendríamos que $q$ divide a $N - p_1 p_2 \cdots p_n = 1$, lo cual es imposible.

Por lo tanto, $q$ no está en la lista $\{p_1, p_2, \ldots, p_n\}$, lo que contradice la suposición de que esta lista contenía todos los primos.

Concluimos que debe existir infinitos números primos.
\end{proof}

\subsection*{Test 18: Irracionalidad de $\sqrt{2}$}

\begin{theorem}
El número $\sqrt{2}$ es irracional, es decir, $\sqrt{2} \notin \mathbb{Q}$.
\end{theorem}

\begin{proof}
Supongamos, por contradicción, que $\sqrt{2}$ es racional. Entonces existen enteros $a, b$ con $b \neq 0$ tales que $\sqrt{2} = \frac{a}{b}$, donde podemos suponer que $\frac{a}{b}$ está en forma reducida, es decir, $\gcd(a, b) = 1$.

Elevando al cuadrado:
\[
2 = \frac{a^2}{b^2} \implies a^2 = 2b^2.
\]

Esto implica que $a^2$ es par. Probemos que entonces $a$ es par. Si $a$ fuera impar, entonces $a = 2k+1$ para algún entero $k$, y:
\[
a^2 = (2k+1)^2 = 4k^2 + 4k + 1 = 2(2k^2 + 2k) + 1,
\]
lo cual es impar, contradiciendo que $a^2$ es par. Por lo tanto, $a$ es par, y podemos escribir $a = 2m$ para algún entero $m$.

Sustituyendo en $a^2 = 2b^2$:
\[
(2m)^2 = 2b^2 \implies 4m^2 = 2b^2 \implies b^2 = 2m^2.
\]

Por el mismo argumento, $b^2$ es par, luego $b$ es par. Pero entonces tanto $a$ como $b$ son pares, lo que contradice la hipótesis de que $\gcd(a, b) = 1$.

Por lo tanto, nuestra suposición inicial es falsa, y $\sqrt{2}$ es irracional.
\end{proof}

\subsection*{Test 19: Paridad y cuadrados (par)}

\begin{proposition}
Si $a^2$ es par, entonces $a$ es par.
\end{proposition}

\begin{proof}
Lo demostramos por contrapositiva. El contrapositivo de ``si $a^2$ es par entonces $a$ es par'' es ``si $a$ es impar entonces $a^2$ es impar''.

Supongamos que $a$ es impar. Entonces $a = 2k + 1$ para algún entero $k$. Entonces:
\[
a^2 = (2k+1)^2 = 4k^2 + 4k + 1 = 2(2k^2 + 2k) + 1.
\]

Como $2k^2 + 2k$ es un entero, $a^2$ tiene la forma $2m + 1$ con $m = 2k^2 + 2k \in \mathbb{Z}$, lo que significa que $a^2$ es impar.

Por lo tanto, por contrapositiva, si $a^2$ es par entonces $a$ es par.
\end{proof}

\subsection*{Test 20: Paridad y cuadrados (impar)}

\begin{proposition}
Si $n^2$ es impar, entonces $n$ es impar.
\end{proposition}

\begin{proof}
Lo demostramos por contrapositiva. El contrapositivo de ``si $n^2$ es impar entonces $n$ es impar'' es ``si $n$ es par entonces $n^2$ es par''.

Supongamos que $n$ es par. Entonces $n = 2k$ para algún entero $k$. Entonces:
\[
n^2 = (2k)^2 = 4k^2 = 2(2k^2).
\]

Como $2k^2$ es un entero, $n^2$ tiene la forma $2m$ con $m = 2k^2 \in \mathbb{Z}$, lo que significa que $n^2$ es par.

Por lo tanto, por contrapositiva, si $n^2$ es impar entonces $n$ es impar.
\end{proof}

\end{document}